% =====================================================================
%  BÁO CÁO TIẾN ĐỘ — BẢN NGẮN (10 slide)
%  Biên dịch:  xelatex baocao_ngan.tex
%  Bản đầy đủ (17 slide + phụ lục): baocao_noidung.tex / .pdf
% =====================================================================
\documentclass[11pt,a4paper]{article}

\usepackage{fontspec}
\setmainfont{Times New Roman}
\usepackage{amsmath}
\usepackage{amssymb}
\usepackage{graphicx}
\usepackage{geometry}
\geometry{margin=2.4cm}
\usepackage{xcolor}
\usepackage{enumitem}
\setlist{nosep,leftmargin=1.3em}
\usepackage{titlesec}
\titleformat{\section}{\normalfont\large\bfseries\color{blue!55!black}}{Slide \thesection.}{0.5em}{}
\titleformat{\subsection}[runin]{\normalfont\bfseries}{}{0em}{}[.]
\titlespacing{\section}{0pt}{14pt}{6pt}

\newcommand{\slide}[1]{\section{#1}}
\newcommand{\noidung}{\subsection*{Nội dung slide}}
\newcommand{\loinoi}{\subsection*{Lời nói}}

\title{\textbf{BÁO CÁO TIẾN ĐỘ DỰ ÁN NGHIÊN CỨU --- BẢN NGẮN}\\ \large
Contagion: đo lan truyền prompt injection trong mạng LLM agent}
\author{Người báo cáo: \underline{\hspace{5cm}} \quad Lớp/Đơn vị: \underline{\hspace{4cm}}}
\date{Ngày báo cáo: \underline{\hspace{4cm}}}

\begin{document}
\maketitle
\thispagestyle{empty}

\vspace{-0.4cm}
\begin{center}
\fbox{\parbox{0.92\textwidth}{
\textbf{Bản ngắn 10 slide.} Mỗi mục gồm \textbf{[Nội dung slide]} để chép lên slide
chiếu, và \textbf{[Lời nói]} là văn nói đầy đủ. Muốn bản chi tiết (17 slide + phụ lục
gồm câu hỏi hội đồng và từ điển thuật ngữ) thì dùng \texttt{baocao\_noidung.pdf}.
}}
\end{center}

\newpage

% =====================================================================
\slide{Bối cảnh và câu hỏi nghiên cứu}

\noidung
\begin{itemize}
  \item LLM agent được triển khai như một \emph{tổ chức}: planner $\to$ worker $\to$
        reviewer $\to$ aggregator. Mỗi lần chuyển giao là một \textbf{ranh giới tin cậy}.
  \item Kẻ tấn công chèn chỉ dẫn độc hại vào nguồn dữ liệu agent đọc; đầu ra của agent
        bị nhiễm trở thành đầu vào của agent sau.
  \item Benchmark hiện có (InjecAgent, AgentDojo) chỉ đo \emph{một agent, một bước}.
  \item \textbf{Câu hỏi:} một agent bị nhiễm thì lan được bao xa, và cái gì quyết định
        nó sống sót qua mỗi lần chuyển giao?
  \item \textbf{Cách làm:} khung đo với \emph{hai protocol độc lập}; 5 model / 5 nhà phát triển;
        3 topology.
\end{itemize}

\loinoi
Kính thưa thầy/cô. Dự án của em nghiên cứu \emph{lan truyền prompt injection trong mạng
nhiều LLM agent}.

Các hệ thống LLM hiện nay không còn là một model đơn lẻ mà là một tổ chức: một agent
lập kế hoạch, chia việc cho agent con, một agent kiểm tra, một agent tổng hợp. Dữ liệu
truyền qua lại là văn bản tự do, nên mỗi lần chuyển giao là một ranh giới tin cậy ---
agent nhận không phân biệt được đâu là dữ liệu, đâu là mệnh lệnh. Kẻ tấn công khai thác
đúng điểm đó: chỉ cần đặt chỉ dẫn độc hại vào một tài liệu hay một kết quả gọi công cụ
mà agent sẽ đọc.

Các benchmark hiện có đều đo \emph{một} agent trong \emph{một} bước, nên không trả lời
được câu hỏi em đặt ra: một khi đã lọt vào, nó lan được bao xa, và cái gì quyết định nó
sống sót qua từng lần chuyển giao. Để trả lời, em xây một khung đo gồm hai protocol độc
lập, chạy trên năm model thuộc năm nhà phát triển độc lập.

% =====================================================================
\slide{Phương pháp: hai protocol độc lập}

\noidung
\begin{itemize}
  \item \textbf{Protocol A --- đo từng cạnh có kiểm soát:} ép đầu ra của agent nguồn
        thành ``đã nhiễm'' $\to$ đưa cho agent nhận $\to$ chấm. Cho
        $s_i=\Pr[C_i{=}1 \mid C_{i-1}{=}1]$.
  \item \textbf{Protocol B --- chạy tự nhiên:} chạy cả mạng, đo tỉ lệ thành công đầu-cuối
        (ASR) và chỉ số lây nhiễm $R_0$.
  \item \textbf{Vì sao hai?} Một cách đo thì khi hai con số lệch nhau, không phân biệt
        được ``đo sai'' với ``hiểu sai hệ thống''.
  \item Chấm điểm \textbf{tất định}, \emph{không} dùng LLM làm giám khảo (tránh giám
        khảo cũng bị lây). 3 ngữ cảnh công việc hợp lệ xoay vòng.
  \item 5 model / 5 nhà phát triển độc lập: Llama 3.3 70B (Meta), DeepSeek V3.2
        (DeepSeek), Claude Sonnet 4.5 (Anthropic), Nova Pro (Amazon), qwen2.5:7b
        (Alibaba, chạy local).
\end{itemize}

\loinoi
Điểm em muốn nhấn mạnh là em cố tình xây \emph{hai} cách đo và giữ chúng tách biệt.

Cách thứ nhất: em không chờ mạng tự lan, mà chủ động \emph{ép} đầu ra của agent trước
thành nội dung đã bị nhiễm, đưa đúng nội dung đó cho agent sau, rồi chấm. Làm nhiều lần
thì được xác suất có điều kiện: nếu agent trước đã nhiễm thì agent sau nhiễm với xác
suất bao nhiêu. Cách thứ hai là để mạng chạy tự nhiên rồi đo tỉ lệ thành công đầu-cuối.

Lý do phải có cả hai: nếu chỉ có một cách đo, khi thấy hai con số không khớp, em không
phân biệt được là do \emph{em đo sai} hay do \emph{hệ thống thực sự khác điều em tưởng}.
Có hai cách độc lập thì mâu thuẫn giữa chúng trở thành một phát hiện.

Việc chấm ``agent có bị nhiễm hay không'' làm bằng công cụ tất định, không dùng LLM làm
giám khảo --- vì trong bài toán lan truyền, giám khảo đọc chính nội dung cần chấm nên có
thể bị ảnh hưởng. Em cũng dùng ba ngữ cảnh công việc hợp lệ khác nhau để kết quả không
phụ thuộc vào một đề bài duy nhất.

% =====================================================================
\slide{Quy mô thực nghiệm: đã chạy những gì và bao nhiêu lần}

\noidung
\begin{itemize}
  \item \textbf{Tổng khối lượng:} 46 thư mục kết quả (25 có cấu trúc chuẩn);
        \textbf{$\approx$ 9.070 lượt gọi model}; \textbf{2.342 trial tự nhiên};
        \textbf{4,3 giờ} thời gian máy ghi lại được.
  \item \textbf{5 model / 5 nhà phát triển} · \textbf{2 protocol độc lập} ·
        \textbf{3 topology} · \textbf{3 defense} · \textbf{3 kiểu biến hình tấn công}.
  \item \textbf{Cấu hình từng loại thí nghiệm:}
        \begin{center}
        \small
        \begin{tabular}{|l|l|}
        \hline
        \textbf{Loại thí nghiệm} & \textbf{Số lần chạy}\\
        \hline
        Ô chuẩn (mọi model) & chuỗi $n=4$: \textbf{40} trial tự nhiên $+$ \textbf{30} trial/cạnh\\
        Đường cong chiều sâu & \textbf{60} trial tự nhiên; \textbf{30}/cạnh (qwen 25);\\
         & chuỗi $n=15$ (Llama), $n=8$ (DeepSeek), $n=7$ (qwen)\\
        Mô hình có vòng & \textbf{40} trial $\times$ \textbf{5 vòng} $\times$ 3 agent\\
        Utility (công việc hợp lệ) & \textbf{20} cặp sạch/bị tấn công cho mỗi defense\\
        Độ nhạy (seed $\times$ temp) & \textbf{8} seed $\times$ \textbf{20} trial/cạnh ở temp $0{,}7$\\
        Biến hình $\times$ defense & \textbf{6} ô $\times$ \textbf{30} mẫu/ô\\
        Topology & 3 topology $\times$ $n=7$: 40 trial tự nhiên $+$ 30/cạnh\\
        \hline
        \end{tabular}
        \end{center}
  \item Cách đếm: lượt gọi là \textbf{ước lượng cận dưới} tính từ cấu hình
        (trial tự nhiên $\times$ số agent $+$ trial/cạnh $\times$ số cạnh), \emph{chưa}
        tính các lần thử lại khi ép đầu ra bị nhiễm và các lần chạy hỏng.
\end{itemize}

\loinoi
Thầy/cô thường hỏi khối lượng thực nghiệm, nên em chuẩn bị riêng một slide cho việc này.

Tính đến nay em đã chạy khoảng \emph{chín nghìn lượt gọi model}, trên \emph{hai nghìn ba
trăm trial tự nhiên}, trải qua bốn mươi sáu thư mục kết quả. Con số chín nghìn là
\emph{ước lượng cận dưới}: em đếm theo công thức từ cấu hình của từng thí nghiệm chứ
không đếm từng lần gọi, và chưa tính các lần thử lại khi phải ép đầu ra thành nội dung bị
nhiễm. Em nói rõ điều này để không ai hiểu nhầm đó là con số đo trực tiếp.

Về cấu hình cụ thể. Mỗi ô chuẩn --- tức là mỗi model, mỗi topology --- chạy bốn mươi lần
tự nhiên, cộng ba mươi lần cho \emph{mỗi cạnh} theo giao thức có kiểm soát. Thí nghiệm
đường cong chiều sâu chạy sáu mươi lần tự nhiên trên chuỗi dài, tới mười lăm agent với
Llama và bảy đến tám agent với hai model còn lại. Thí nghiệm mô hình có vòng chạy bốn
mươi lần, mỗi lần năm vòng. Phần đo ảnh hưởng lên công việc hợp lệ chạy hai mươi cặp
sạch và bị tấn công cho mỗi cơ chế phòng thủ. Phần độ nhạy chạy tám seed, mỗi seed hai
mươi lần mỗi cạnh. Bảng biến hình tấn công có sáu ô, mỗi ô ba mươi mẫu.

Có hai điểm em muốn nói thêm về con số này. Thứ nhất, thời gian máy ghi lại được là bốn
phẩy ba giờ, nhưng đây cũng là cận dưới, vì hai họ script --- đường cong chiều sâu và mô
hình có vòng --- không ghi lại thời gian chạy. Thứ hai, điều quan trọng không phải là con
số chín nghìn, mà là \emph{mỗi ô đều có khoảng tin cậy và đều công bố độ lệch nhỏ nhất
phát hiện được}. Chạy nhiều mà không biết mình phát hiện được gì thì con số lớn cũng
không có ý nghĩa.

% =====================================================================
\slide{Kết quả 1 và 2: ``giả định Markov'' là đẳng thức, và nó \emph{thất bại}}

\noidung
\begin{itemize}
  \item Trong chuỗi: $C_i{=}1 \Rightarrow C_{i-1}{=}1$
        $\Rightarrow \mathrm{ASR}=\prod_i s_i^{\mathrm{nat}}$ là \textbf{đẳng thức}, không
        phải giả định (kiểm: chuỗi 4 agent $0{,}850=0{,}850$; chuỗi 7 agent
        $0{,}450=0{,}450$).
  \item Vậy câu hỏi thật là: con số đo \emph{cách ly} có bằng con số \emph{trong chuỗi}?
        Câu trả lời: \textbf{tuỳ model}.
        \begin{itemize}
          \item DeepSeek: khớp tới $2{,}2\%$ $\Rightarrow$ chuyển được.
          \item Llama: cạnh giữa cách ly $0{,}233$ nhưng trong chuỗi $0{,}875$
                $\Rightarrow$ \textbf{lệch $3{,}75$ lần}.
        \end{itemize}
  \item \textbf{Sai số lớn dần theo độ sâu} --- đo trên chuỗi dài, 3 model:
        \begin{center}
        \small
        \begin{tabular}{|l|c|c|c|c|}
        \hline
        \textbf{Model} & \textbf{Chuỗi} & \textbf{Sai số} & \textbf{$\rho$} & \textbf{Dốc (\%/hop)}\\
        \hline
        qwen2.5:7b & n=7 & $31\%\to 89\%$ & $+1{,}00$ & $+11{,}8$\\
        Llama 3.3 70B & n=15 & $0\%\to 92\%$ & $+0{,}77$ & $+6{,}8$\\
        DeepSeek V3.2 & n=8 & $7\%$--$35\%$ & $-0{,}07$ & $+0{,}1$\\
        \hline
        \end{tabular}
        \end{center}
  \item $\Rightarrow$ \textbf{độ dốc đường cong là phép chẩn đoán}: một chuỗi dài, một
        lần chạy, biết ngay model đó có cộng dồn được hay không.
\end{itemize}

\loinoi
Đây là kết quả trung tâm. Trước hết, một kết quả lý thuyết: trong chuỗi, để agent thứ ba
bị nhiễm thì agent thứ hai phải đã nhiễm, vì chính đầu ra của agent thứ hai mới lây sang
agent thứ ba. Hệ quả là tỉ lệ thành công đầu-cuối \emph{đúng bằng} tích các xác suất có
điều kiện --- một đẳng thức, không phải giả định cần kiểm định. Em kiểm bằng số đo thật:
chuỗi bốn agent cho $0{,}850$ và tích cũng $0{,}850$; chuỗi bảy agent cho $0{,}450$ và
tích cũng $0{,}450$.

Vậy câu hỏi ``Markov có đúng không'' là câu hỏi đặt sai. Câu hỏi đúng là: con số đo
\emph{cách ly} --- đo riêng lẻ từng cạnh --- có bằng con số khi mạng chạy thật không.
Câu trả lời là \emph{tuỳ model}. Với DeepSeek thì khớp tới hai phần trăm. Với Llama thì
ngược lại: một cạnh ở giữa đo cách ly cho $0{,}233$ nhưng trong chuỗi thật là $0{,}875$
--- đo cách ly đánh giá thấp gần bốn lần.

Đáng chú ý hơn là \emph{quy luật} của sai số. Em chạy một chuỗi dài, mỗi nút là một mốc
độ sâu, rồi so tích các con số cách ly với xác suất nhiễm thật. Với qwen, sai số tăng
đơn điệu từng bước, tương quan hạng đúng bằng $+1$. Với Llama, tương quan $+0{,}77$, lên
tới $92\%$. Với DeepSeek thì không có xu hướng nào.

Ý nghĩa thực tiễn: độ dốc của đường cong tự nó là phép chẩn đoán. Chỉ cần một chuỗi dài,
một lần chạy, nhìn độ dốc là biết có nên tin vào việc nhân các con số đo cách ly hay
không. Một chi tiết về trung thực: bốn điểm cuối của Llama không xác định được vì chuỗi
đã chết hẳn, nên em \emph{loại} chứ không gán cho chúng sai số 100\% cho đẹp hình.

% =====================================================================
\slide{Cơ chế: hình thức của thông điệp quyết định}

\noidung
\begin{itemize}
  \item Cùng mục tiêu, cùng agent nhận, cùng cách chấm, cùng phân công --- chỉ đổi cách
        trình bày:
        \begin{center}
        \small
        \begin{tabular}{|l|c|}
        \hline
        \textbf{Cách trình bày} & \textbf{Xác suất sống sót}\\
        \hline
        Một câu khẳng định trần & $0{,}067$\\
        Nhúng trong một sản phẩm công việc & $\mathbf{0{,}967}$\\
        \hline
        \end{tabular}
        \end{center}
  \item Chênh \textbf{14 lần} chỉ vì cách đóng gói.
  \item Đây là lời giải thích cho kết quả trước: artefact ``chuẩn hoá'' dùng khi đo cách
        ly \emph{không đại diện} cho thông điệp thật trong mạng.
  \item Tự kiểm tra lại: phát lại nội dung bắt được khớp \textbf{2/3} cạnh --- báo cáo
        đúng $2/3$, không làm tròn thành ``khớp hoàn toàn''.
\end{itemize}

\loinoi
Kết quả trước nói rằng phép đo cách ly sai. Câu hỏi tiếp là sai vì cái gì. Câu trả lời:
vì \emph{hình thức} của thông điệp.

Em làm một thí nghiệm đối chứng rất gọn: giữ nguyên mọi thứ --- cùng mục tiêu, cùng các
agent nhận, cùng cách chấm, cùng phân công --- chỉ đổi cách trình bày. Nếu nội dung độc
hại đến dưới dạng một câu khẳng định trần, tỉ lệ sống sót là $0{,}067$, tức gần như bị
chặn hết. Nếu nó được nhúng vào trong một sản phẩm công việc thật, tỉ lệ sống sót là
$0{,}967$. Chênh nhau mười bốn lần, chỉ vì cách đóng gói.

Điều này giải thích vì sao phép đo cách ly lệch: khi đo cách ly, em buộc phải dùng một
artefact ``chuẩn hoá'' để đưa cho agent nhận, và artefact đó không đại diện cho thông
điệp thật trong mạng.

Em cũng tự kiểm tra lại chính mình bằng cách phát lại nội dung đã bắt được trong chuỗi
thật. Kết quả là hai trên ba cạnh tái hiện được, và em báo cáo đúng hai trên ba.

% =====================================================================
\slide{Kết quả 3: tiêu chí ngưỡng dịch tễ \emph{vô hiệu} với mạng agent}

\noidung
\begin{itemize}
  \item Tiêu chí quen thuộc: ``$R_0<1$ thì dịch tắt, hệ an toàn''.
  \item \textbf{Định lý:} ma trận lây nhiễm của \emph{mọi} mạng không có vòng, sắp theo
        topology, là \textbf{tam giác chặt} $\Rightarrow$ mọi giá trị riêng bằng $0$
        $\Rightarrow \rho(M)=0$ \textbf{luôn luôn}, bất kể cạnh mạnh cỡ nào.
  \item Nghĩa là tiêu chí được thoả mãn \emph{tầm thường}: không phân biệt được hệ an
        toàn với hệ nguy hiểm.
  \item Bằng chứng thực nghiệm:
        \begin{itemize}
          \item chuỗi 7 agent: $R_0=0{,}821<1$ nhưng ASR $=0{,}450$
          \item hình sao: $R_0=0{,}420<1$ nhưng ASR $=0{,}725$
        \end{itemize}
  \item \textbf{Độ lan tới} và \textbf{khả năng tái lây} là hai rủi ro khác nhau.
\end{itemize}

\loinoi
Kết quả thứ ba là phản biện một công cụ đang được dùng rộng rãi. Trong mạng lưới nói
chung, người ta hay dùng tiêu chí: nếu chỉ số lây nhiễm nhỏ hơn một thì dịch sẽ tắt.

Em chứng minh rằng với mạng agent dạng feed-forward --- mạng không có vòng, mọi agent
chỉ chuyển tiếp về phía trước --- thì tiêu chí này \emph{luôn luôn} được thoả mãn một
cách tầm thường. Lý do rất gọn: ma trận mô tả ``một agent lây cho bao nhiêu agent'', khi
sắp theo thứ tự truyền, luôn có dạng tam giác; mà ma trận tam giác thì toàn bộ giá trị
riêng bằng không. Không phụ thuộc cạnh mạnh hay yếu. Tức là tiêu chí đó luôn báo ``an
toàn''.

Thực nghiệm xác nhận: chuỗi bảy agent có chỉ số lây nhiễm $0{,}821$, nhỏ hơn một, nhưng
vẫn có $45$ phần trăm số lần tấn công thành công đầu-cuối. Hình sao còn rõ hơn: chỉ số
$0{,}420$ mà tỉ lệ thành công tới $72{,}5$ phần trăm. Nghĩa là ``đi được tới đích'' và
``tái lây được tiếp'' là hai rủi ro khác nhau.

% =====================================================================
\slide{Kết quả 3 (tiếp): có vòng lặp thì tiêu chí \emph{có} hiệu lực --- và thành quy tắc thiết kế}

\noidung
\begin{itemize}
  \item Mô hình thật: manager giao việc cho $w$ worker, worker \textbf{báo cáo ngược lại}.
  \item \textbf{Công thức} (kiểm bằng số tới $4{,}4\times10^{-16}$ trên 320 cấu hình):
        \[
          \rho(M)=\sqrt{\textstyle\sum_j a_jc_j}\quad\text{với \emph{mọi} } w,
          \qquad\text{vượt ngưỡng}\iff \sum_j a_jc_j>1 .
        \]
        Không có báo cáo ngược ($c_j{=}0$) $\Rightarrow \rho=0$ bất kể cạnh mạnh cỡ nào.
  \item Đối xứng ($a_j{=}c_j{=}s$): $w\,s^2>1$ --- ở $s{=}0{,}5$ cần \textbf{5} worker,
        $0{,}7$ cần \textbf{3}, $0{,}9$ cần \textbf{2}.
  \item \textbf{Hai dự đoán ngược nhau, cả hai đều đúng:}
        \begin{center}
        \small
        \begin{tabular}{|l|c|c|l|}
        \hline
        \textbf{Model} & $\sum_j a_jc_j$ & $\rho$ & \textbf{Dự đoán $\to$ Đo được}\\
        \hline
        Llama & $1{,}800$ & $\mathbf{1{,}342}$ & bền $\to$ manager $0{,}725$, vòng 5 còn $0{,}700$ \checkmark\\
        qwen & $0{,}578$ & $0{,}760$ & yếu dần $\to$ manager $0{,}775\to\mathbf{0{,}650}$ \checkmark\\
        \hline
        \end{tabular}
        \end{center}
  \item Giới hạn đã ghi trong bài: mới 5 vòng --- đủ phân biệt ``yếu dần'' với ``bão
        hoà'', \emph{chưa} đủ chứng minh tắt hẳn.
\end{itemize}

\loinoi
Nhưng tiêu chí ngưỡng không phải lúc nào cũng vô nghĩa --- nó chỉ vô nghĩa khi mạng
không có vòng. Mà mạng agent thật thì \emph{có} vòng: mô hình rất phổ biến là một manager
giao việc cho các worker, rồi worker báo cáo kết quả ngược trở lại. Chính các cạnh báo
cáo ngược đó tạo thành vòng.

Em tính được công thức bán kính phổ cho mô hình này: căn bậc hai của tổng, qua từng
worker, của tích ``xác suất lây xuống nhân xác suất báo cáo ngược lên''. Em kiểm công
thức bằng số trên hơn ba trăm cấu hình ngẫu nhiên, sai số lớn nhất là bốn phần mười mũ
mười sáu --- tức là một đẳng thức, không phải xấp xỉ.

Từ đó ra một \emph{quy tắc thiết kế}: hệ bắt đầu tự duy trì lây nhiễm khi tổng đó lớn hơn
một. Trường hợp đối xứng thì điều kiện là $w$ nhân $s$ bình phương lớn hơn một. Cụ thể:
xác suất sống sót mỗi chiều là $0{,}5$ thì cần năm worker mới vượt ngưỡng; $0{,}7$ thì
cần ba; $0{,}9$ thì chỉ cần hai.

Điểm em thấy quan trọng nhất: công thức này \emph{dám} đưa ra hai dự đoán \emph{ngược
nhau} trên hai model em đã đo, và em đã chạy cả hai. Với Llama, tổng bằng $1{,}8$, dự
đoán vượt ngưỡng --- và thực tế manager bị tái nhiễm bởi chính báo cáo của worker mình
với xác suất $0{,}725$, tới vòng năm vẫn còn $0{,}700$, tức hệ không tắt. Với qwen, tổng
chỉ $0{,}578$, dự đoán yếu dần --- và thực tế manager từ $0{,}775$ giảm còn $0{,}650$ ở
vòng năm, worker dừng lại chứ không bão hoà.

Em xin nói rõ giới hạn: mới chạy năm vòng, đủ để phân biệt ``yếu dần'' với ``bão hoà''
nhưng chưa đủ để khẳng định tắt hẳn, nên trong bài em không dùng chữ ``tắt hẳn''.

% =====================================================================
\slide{Kết quả 4: đánh giá phòng thủ cần baseline ``không phòng thủ''}

\noidung
\begin{itemize}
  \item Cùng defense, cùng kiểu tấn công, cùng code chấm điểm, chỉ đổi model:
        \begin{itemize}
          \item Claude: không phòng thủ đã chỉ $0{,}10$; bật redaction còn $0{,}00$
                $\Rightarrow$ ``trông hoàn hảo'' nhưng chỉ có $10\%$ để bảo vệ.
          \item DeepSeek: không phòng thủ $1{,}00$; bật redaction còn $0{,}00$
                $\Rightarrow$ defense thật sự làm hết việc.
        \end{itemize}
  \item \textbf{Kết luận:} thiếu baseline ``không phòng thủ'' thì \emph{không xác định
        được} hiệu quả của một defense.
  \item Kèm theo: biến hình tấn công (tách chuỗi, chèn khoảng trắng) phá được redaction.
  \item Phải báo cáo \emph{cặp} (ASR, mức mất công việc hợp lệ): trên DeepSeek,
        \texttt{paraphrase} \textbf{không giảm} ASR mà còn làm mất thêm công việc hợp lệ
        --- tệ hơn cả không làm gì.
\end{itemize}

\loinoi
Kết quả thứ tư là về \emph{cách đánh giá}. Em dùng cùng một cơ chế phòng thủ, cùng kiểu
tấn công, cùng đoạn code chấm điểm, chỉ đổi model.

Trên Claude, khi \emph{không} bật phòng thủ, tỉ lệ nhiễm đã chỉ là $0{,}10$; bật
redaction lên thì còn $0{,}00$. Nhìn con số cuối thì có vẻ hoàn hảo, nhưng thực ra nó chỉ
có mười phần trăm để bảo vệ --- phần còn lại là model vốn đã kháng tốt. Trên DeepSeek,
cùng phép so sánh lại có ý nghĩa khác hẳn: không phòng thủ là $1{,}00$, bật lên còn
$0{,}00$.

Nên khuyến nghị phương pháp là: báo cáo đánh giá phòng thủ mà không kèm con số baseline
khi không phòng thủ thì người đọc \emph{không thể} biết phòng thủ đó có thật sự làm việc
hay không. Em gọi đây là hiệu ứng sàn.

Thêm một điểm nữa: phải báo cáo cả \emph{công việc hợp lệ còn làm được bao nhiêu}. Trên
DeepSeek, cơ chế viết lại ngữ nghĩa \texttt{paraphrase} không giảm được lan truyền chút
nào mà còn làm mất thêm công việc hợp lệ --- tệ hơn cả việc không làm gì. Nếu chỉ báo cáo
mỗi tỉ lệ tấn công thì nó trông ``trung tính''; chỉ khi báo cáo cặp hai con số mới thấy
vấn đề.

% =====================================================================
\slide{Kết quả phụ: vai trò agent và năng lực model}

\noidung
\begin{itemize}
  \item Cùng model, cùng kiểu tấn công: xác suất sống sót chênh \textbf{$6{,}5$ lần} giữa
        các vai trò (worker $0{,}133$ so với reviewer $0{,}867$).
  \item ``Dễ bị lây'' \textbf{không} đi theo ``model mạnh hơn'':
        \begin{center}
        \small
        \begin{tabular}{|l|c|c|}
        \hline
        \textbf{Model} & \textbf{ASR} & \textbf{$\bar s$}\\
        \hline
        Llama 3.3 70B & $0{,}800$ & $0{,}722$\\
        DeepSeek V3.2 & $0{,}475$ & $0{,}933$\\
        qwen2.5:7b & $0{,}300$ & $0{,}611$\\
        Nova Pro & $0{,}075$ & $0{,}178$\\
        Claude Sonnet 4.5 & $0{,}000$ & $0{,}211$\\
        \hline
        \end{tabular}
        \end{center}
  \item \textbf{Topology} ở cùng 7 agent, cùng $\bar s\approx0{,}76$: ASR $0{,}450$ /
        $0{,}725$ / $0{,}800$ cho chain / star / tree.
  \item \textbf{Đặt phòng thủ ở đâu:} trên mạng không vòng, chỉ agent trên đường tới mục
        tiêu mới có ý nghĩa, và mọi agent trên đường đó \emph{hiệu quả như nhau}
        $\Rightarrow$ bài toán \emph{cấu trúc}, không phải thống kê.
\end{itemize}

\loinoi
Hai kết quả phụ, nhưng theo em là quan trọng cho người thiết kế hệ thống.

Thứ nhất, xác suất lây không phải thuộc tính của model mà là thuộc tính của \emph{vị trí}
agent trong quy trình. Cùng một model và cùng kiểu tấn công, agent ở vai trò worker có
xác suất bị lây $0{,}133$, còn reviewer là $0{,}867$ --- chênh sáu lần rưỡi.

Thứ hai, trực giác ``model mạnh hơn thì an toàn hơn'' không đúng. Llama 3.3 70B có tỉ lệ
tấn công thành công $0{,}800$, cao hơn DeepSeek $0{,}475$, cao hơn qwen $0{,}300$, cao
hơn Nova Pro $0{,}075$, và Claude gần như bằng không. Đây không phải thứ tự năng lực. Nên
phải đo trên đúng model mình dùng, không suy ra từ bảng xếp hạng. Topology cũng quyết
định: ba mạng bảy agent có cùng xác suất lây trung bình nhưng cho tỉ lệ thành công khác
hẳn nhau.

Cuối cùng, từ định lý ở kết quả ba, em rút ra câu trả lời cho câu hỏi ``đặt phòng thủ ở
đâu'': trên mạng không có vòng, chỉ agent nằm trên đường từ điểm xâm nhập tới mục tiêu
mới đáng quan tâm, và trong số đó mọi agent đều hiệu quả như nhau. Nghĩa là đây là bài
toán về cấu trúc đồ thị và chi phí triển khai, không phải bài toán cần số liệu thống kê.

% =====================================================================
\slide{Trung thực phương pháp: một phát hiện của chính em đã bị rút lại}

\noidung
\begin{itemize}
  \item Ban đầu: trên DeepSeek, kiểm định bác bỏ giả thuyết Markov với
        $p=\mathbf{0{,}0030}$ --- tưởng là phát hiện lớn.
  \item Nguyên nhân thật: protocol khi đó \emph{giữ cố định} artefact đã nhiễm và dùng
        lại cho mọi trial $\Rightarrow$ các trial phụ thuộc nhau, độ lệch bị phóng đại.
  \item Sửa: thêm chế độ \texttt{--fresh-artifact} (mỗi trial rút artefact mới)
        $\Rightarrow$ chạy lại: $p=\mathbf{0{,}630}$, \emph{không} bác bỏ được gì.
  \item Giữ cả cái sai và cách sửa trong bài: với nghiên cứu đo lường,
        \textbf{cách đo là một phần của kết quả}.
  \item Ngược lại, bất thường của Llama \emph{đứng vững} qua cả hai chế độ artefact.
\end{itemize}

\loinoi
Phần này em dành riêng để nói về một sai sót của chính mình, vì em nghĩ nó là một phần
quan trọng của tiến độ.

Giai đoạn đầu, em chạy kiểm định thống kê và thấy trên DeepSeek giả thuyết Markov bị bác
bỏ với $p$ bằng $0{,}003$ --- một kết quả rất mạnh, đủ để viết thành đóng góp chính.
Nhưng khi kiểm tra lại quy trình đo, em phát hiện protocol khi đó \emph{giữ cố định} một
artefact đã bị nhiễm và dùng lại cho mọi trial. Cách làm đó khiến các trial không còn độc
lập và làm độ lệch bị phóng đại.

Em sửa bằng cách rút artefact mới cho từng trial rồi chạy lại: $p$ bằng $0{,}630$, hoàn
toàn không bác bỏ được gì. Kết quả cũ là artefact của thiết kế đo.

Em giữ cả hai trong bài --- cả cái sai và cách sửa --- vì trong nghiên cứu đo lường, cách
đo là một phần của kết quả. Để công bằng, em cũng kiểm ngược: kết quả bất thường của
Llama thì đứng vững qua cả hai chế độ, tức đó là bất thường thật.

% =====================================================================
\slide{Sản phẩm, hạn chế, kế hoạch}

\noidung
\begin{itemize}
  \item \textbf{Sản phẩm:} bản thảo bài báo định dạng AAMAS 2027 --- \textbf{8 trang nội
        dung} + 1 trang tài liệu tham khảo (đúng giới hạn); 8 bảng + 2 hình; mọi số sinh
        tự động từ file kết quả, \emph{không nhập tay}.
  \item Kèm theo: tài liệu tái lập (mỗi bảng gắn với câu lệnh sinh ra nó) và script đối
        chiếu từng con số trong bài với kết quả gốc.
  \item \textbf{Hạn chế (đã ghi rõ trong bài):} utility đo bằng proxy; cỡ mẫu $n=30$--$40$
        (chỉ phát hiện được lệch $>0{,}26$); chưa có topology lưới/tranh luận; mô hình có
        vòng mới có 1 instance; transportability mới chứng minh \emph{tồn tại} thất bại;
        chưa hiệu chỉnh đa so sánh.
  \item \textbf{Kế hoạch:} mở rộng topology (lưới, tranh luận); nhiệm vụ có đáp án kiểm
        chứng được để đo utility thật; chạy lại các ô quan trọng ở cỡ mẫu lớn hơn.
\end{itemize}

\loinoi
Về sản phẩm, hiện em có một bản thảo bài báo hoàn chỉnh theo định dạng hội nghị AAMAS
2027, dài đúng tám trang nội dung cộng một trang tài liệu tham khảo. Nguyên tắc em đặt ra
từ đầu là \emph{không có con số nào nhập tay}: mọi bảng hình đều sinh tự động từ file kết
quả, và em có thêm một script đối chiếu từng con số trong bài với kết quả gốc --- nó đã
phát hiện và giúp em sửa hai chỗ ghi thiếu chính xác.

Về hạn chế, em xin nói thẳng. Đại lượng đo mức ảnh hưởng lên công việc hợp lệ hiện là
đại lượng thay thế, chưa phải độ đúng thật. Cỡ mẫu phần lớn là ba mươi đến bốn mươi lần
chạy, ở cỡ đó chỉ phát hiện được lệch lớn hơn $0{,}26$, và em công bố con số này trong
bài chứ không giấu. Em chưa có topology mạng lưới và mạng tranh luận; mô hình có vòng mới
có một instance; và kết quả về tính chuyển được mới chứng minh sự \emph{tồn tại} của thất
bại, chưa ước lượng được tần suất.

Kế hoạch tiếp theo là mở rộng độ phủ topology, đặc biệt là mạng có đường song song vì đó
là nơi bài toán đặt phòng thủ trở thành bài toán tối ưu thật sự; và xây một họ nhiệm vụ
có đáp án kiểm chứng được để đo mức ảnh hưởng lên công việc một cách chính xác.

% =====================================================================
\slide{Kết luận}

\noidung
\begin{itemize}
  \item Lan truyền prompt injection trong mạng agent là hiện tượng \textbf{đa tác nhân}:
        tầm với của một vụ nhiễm là thuộc tính của \emph{đồ thị, vai trò và thông điệp}.
  \item \textbf{Ba thông điệp:}
        \begin{enumerate}
          \item Con số ASR của benchmark một-agent \textbf{không} dùng được như tham số
                để nhân lên thành dự đoán cho pipeline; sai số \emph{tăng theo độ dài
                pipeline} (và độ dốc đó là phép chẩn đoán).
          \item ``$R_0<1$ là an toàn'' \textbf{không} dùng được cho mạng không vòng; khi
                có vòng thì tiêu chí \emph{có} hiệu lực, và trở thành quy tắc thiết kế
                ($\sum_j a_jc_j>1$) đã được kiểm hai chiều trên hai model.
          \item Đánh giá defense mà thiếu baseline ``không phòng thủ'' thì
                \textbf{không xác định được} hiệu quả.
        \end{enumerate}
  \item Sản phẩm: khung đo tái lập được + bản thảo 8 trang + toàn bộ dữ liệu thô, kể cả
        một phát hiện sai của chính nhóm đã được rút lại.
\end{itemize}

\loinoi
Em xin kết luận. Điểm cốt lõi: lan truyền prompt injection trong mạng agent là hiện tượng
đa tác nhân, và tầm với của một vụ nhiễm là thuộc tính của đồ thị, của vai trò agent và
của hình thức thông điệp --- không phải thuộc tính của một agent đơn lẻ.

Ba thông điệp em muốn mang về. Một: con số tỉ lệ tấn công thành công mà các benchmark đo
trên một agent không dùng được như tham số để nhân lên thành dự đoán cho cả pipeline, và
mức sai lệch tăng theo độ dài pipeline --- hơn nữa, độ dốc của đường cong đó tự nó là
phép chẩn đoán. Hai: tiêu chí ngưỡng ``chỉ số lây nhiễm nhỏ hơn một thì an toàn'' không
dùng được cho mạng không có vòng, nhưng khi mạng có vòng thì tiêu chí trở nên có hiệu
lực và trở thành một quy tắc thiết kế, đã được kiểm hai chiều. Ba: đánh giá một cơ chế
phòng thủ mà thiếu con số baseline khi không phòng thủ thì không thể xác định được hiệu
quả.

Sản phẩm bàn giao gồm khung đo tái lập được, bản thảo tám trang, và toàn bộ dữ liệu thô
--- trong đó em giữ lại cả một phát hiện sai của chính mình và cách sửa nó.

Em xin hết phần trình bày và sẵn sàng nhận câu hỏi của thầy/cô.

\end{document}
